\scanpage{49}
\diagram{\begin{tikzcd}[column sep=small,row sep=large]
\text{Fibré des tores maximaux de }G\text{ est} & P/N_1 \simeq Q/N'_1 & {}=\mathcal C(G)\\
\text{--- Borel ---} & P/B_1\simeq Q/B'_1 & {}=\mathcal D(G)\\
\text{--- couples de Killing ---}\arrow[d,equal] & P/T_1\simeq Q/T'_1\arrow[d,equal] & {}=\mathcal C'(G)\arrow[u]\arrow[uu,bend right=50,"\text{fibré principal sous }W_1"']\\
\text{--- }\left\{\begin{array}{l}\text{Tores maximaux \unsure{piqués|pointés} de }G\\\text{Immersions intérieures }T_1\to G\end{array}\right\} & P/T_1 &
\end{tikzcd}}
\margin{fibré \deleted{\unsure{en fibres}} \unsure{sous} $B_1/T_1=\mathcal C(B_1)$ associé à un fibré principal \unsure{homogène} de gr $B_1$}

\uline{N.B.} Soit $T$ un tore \inserted{maximal de $G$}, il y a correspondance biunivoque (\unsure{fonctorielle} $S'/S$) entre les ensembles suivants

a) L'ens des Borel contenant $T$

b) L'ens des isom $T_1-T$ qui sont intérieurs
\margin{$\left\}\unsure{P/T_1}\simeq P'/T'_1\right.$}

\deleted{\unsure{Puis qui}}\unsure{plus}, \deleted{\unsure{on a les conditions}} \inserted{la \unsure{définition} est \unsure{généralisée} par b) :} \deleted{il existe une \unsure{application}}

b') L'ens des isom $T_1\longrightarrow T$ qui sont ``géom$^{\mathrm t}$ \unsure{extérieurs|intérieurs}''
\margin{$\left\}P'/T_1\right.$}
\deleted{\unsure{Autrement dit}} Il existe une projection de b') sur b).

\margin{$W_1$ et $W$ sont \unsure{commutants} l'un de l'autre \unsure{dans cet ens} $P_1$ \unsure{supposé} \unsure{principal} \unsure{lorsque les deux schémas sont} $\ne\phi$]}

De plus les conditions suivantes sur $T$ sont équivalentes

$\alpha$) $T$ est contenu dans un Borel

$\beta$) Il existe un isom intérieur $T_1-T$

$\gamma$) Il existe un isom \unsure{extérieur} $T_1-T$

On dit alors que $T$ est un ``Tore de \unsure{Killing}''. Pour qu'il existe, il f. et s. que $G$ soit isomorphe à une forme extérieure type.

L'image inverse de $\xi\in H^1(S,G_1)$ dans $H^1(S,B_1)$ est en correspondance biunivoque avec l'ens des sous-groupes de Borel de $G$, modulo automorphismes intérieurs ; Il est non vide [i.e. on peut restreindre le groupe structural à $B_1$] dès qu'il existe un sous-groupe de Borel; dans ce cas, il est réduit à un élément : ces deux sous-groupes de Borel de $G$ sont tous conjugués. \deleted{\unsure{Cela est}}

\margin{On n'a \unsure{besoin} que $G_1$ soit une forme extérieure type, mais seulement qu'il y ait un ss-groupe de Borel}

[\uline{Question} ? Cette \deleted{\unsure{condition est vérifiée}} \inserted{\unsure{dernière} condition} \unsure{est-elle toujours} ? Si $S$ est local ?]
\margin{$\times$}
\margin{48}
